Fix \(n\) and \(P\in\mathcal P_{\mathrm G,n}^{\mathrm{JMS}}\).  We first prove the pointwise assertion.  Because \(|Y|\le C_q\) almost surely, conditional Jensen's inequality gives
\[
 \left|E_P[Y\mid X,T]\right|
 \le E_P[|Y|\mid X,T]
 \le C_q
 \quad\text{almost surely}. \tag{1}
\]
Likewise, the definition \(q_0(X)=E_P[Y\mid X]\) implies
\[
 |q_0(X)|\le C_q
 \quad\text{almost surely}. \tag{2}
\]
By the displayed definition
\(
 \xi=Y-q_0(X)-\theta_0\eta
\)
and the class restriction \(E_P[\xi\mid X,T]=0\), while \(q_0(X)\) and \(\eta=T-g_0(X)\) are measurable with respect to \((X,T)\),
\[
 E_P[Y\mid X,T]=q_0(X)+\theta_0\eta
 \quad\text{almost surely}. \tag{3}
\]
Suppose for contradiction that \(\theta_0\ne0\).  Under the Gaussian-class definition, \(\eta\perp X\) and \(\eta\sim N(0,\sigma^2)\), with \(\sigma>0\).  Hence, for every real \(u\),
\[
 P\bigl(|u+\theta_0\eta|>C_q\bigr)>0, \tag{4}
\]
because a nondegenerate Gaussian law assigns positive probability to every nonempty open tail.  Apply \((4)\) conditionally with \(u=q_0(X)\), which is finite by \((2)\).  Independence yields
\[
 P\bigl(|q_0(X)+\theta_0\eta|>C_q\mid X\bigr)>0
 \quad\text{almost surely}. \tag{5}
\]
Taking expectations in \((5)\) gives
\[
 P\bigl(|q_0(X)+\theta_0\eta|>C_q\bigr)>0. \tag{6}
\]
Equations \((1)\) and \((3)\) say that the probability in \((6)\) is zero, a contradiction.  Therefore
\[
 \theta_0(P)=0
 \qquad\text{for every }P\in\mathcal P_{\mathrm G,n}^{\mathrm{JMS}}. \tag{7}
\]

Now suppose \(\mathcal P_{\mathrm G,n}^{\mathrm{JMS}}\ne\varnothing\).  The statistic \(\widetilde\theta\equiv0\) is a single measurable sample map, independent of \(P\), and \((7)\) gives
\[
 \sup_{P\in\mathcal P_{\mathrm G,n}^{\mathrm{JMS}}}
 E_P[(\widetilde\theta-\theta_0(P))^2]=0. \tag{8}
\]
Since squared loss is nonnegative, the definition in \(\text{def:minimax-risks}\) and \((8)\) imply
\[
 \mathfrak R_n^{\mathrm G}=0. \tag{9}
\]
For the same statistic the absolute error is identically zero, and the definition of the generalized quantile gives
\[
 Q_{P,1-\gamma}(|\widetilde\theta-\theta_0(P)|)=0
 \quad\text{for every }P\in\mathcal P_{\mathrm G,n}^{\mathrm{JMS}}.
\]
Nonnegativity and the minimax definition therefore yield
\[
 \mathfrak M_{n,1-\gamma}^{\mathrm G}=0. \tag{10}
\]
In particular, because \(c_\gamma>0\), \(\sigma>0\), and \(n<\infty\), the displayed quantity \(L_{\mathrm G,n,\gamma}\) is strictly positive whenever it is defined.  Thus a claim that the class in \(\text{def:gaussian-class}\) has generalized-quantile risk at least \(L_{\mathrm G,n,\gamma}\) contradicts \((10)\).  The cited Jin--Mackey--Syrgkanis theorem consequently cannot be imported under this simultaneous bounded-outcome and conditional-mean specification.